% Source: https://artofproblemsolving.com/community/c3h1407970_slope_problem
% Stable identity: aops-topic-1407970-q01
\textbf{(a)} A line with $y$-intercept $K$ passes through $(0,K)$. Since $(0,K)\ne(3,3)$, there is exactly one line through $(3,3)$ and $(0,K)$, and different values of $K$ give different lines. The only line through $(3,3)$ with no $y$-intercept is the vertical line $x=3$. So the lines correspond exactly to the integers $K$ with $-3000\le K\le 3000$, and there are $\boxed{6001}$ of them.

\textbf{(b)} For $K\ne3$ the line through $(0,K)$ and $(3,3)$ has slope $m=\dfrac{3-K}{3}$, so its equation is $y=\dfrac{3-K}{3}x+K$ and its $x$-intercept is
\[
x_0=-\dfrac{K}{m}=\dfrac{3K}{K-3}=3+\dfrac{9}{K-3}.
\]
(For $K=3$ the line is $y=3$, which has no $x$-intercept, and $x=3$ has no $y$-intercept.)

\begin{asy}
usepackage("amsmath");
size(6cm);
defaultpen(fontsize(10pt));
import graph;
xaxis("$x$",-1.5,8.5,Arrow(TeXHead)); yaxis("$y$",-1.5,8.5,Arrow(TeXHead));
pair P=(3,3);
draw((-1,7)--(7.5,-1.5),blue+linewidth(0.9));
draw((-1,3)--(8,3),gray+dashed); draw((3,-1)--(3,8),gray+dashed);
dot(P); label("$(3,3)$",P,NE);
dot((0,6),blue); label("$(0,K)$",(0,6),W,blue);
dot((6,0),blue); label("$\left(\dfrac{3K}{K-3},0\right)$",(6,0),2*NE,blue);
label("$y=3$",(8,3),N,gray); label("$x=3$",(3,8),E,gray);
label("\mbox{}",(min(currentpicture,user=true)+max(currentpicture,user=true))/2,black); // keeps the text after the figure black when LaTeX typesets the labels inline
\end{asy}

So $x_0$ is an integer exactly when $K-3$ divides $9$, that is $K-3\in\{\pm1,\pm3,\pm9\}$, giving $K\in\{4,2,6,0,12,-6\}$ with $x$-intercepts $12,-6,6,0,4,2$ respectively. There are $\boxed{6}$ such lines (for $K=0$ the line $y=x$ passes through the origin, and both intercepts are $0$).
